\documentclass{fhe-shared/luxfhe}
\usepackage{fhe-shared/luxcover}
\input{fhe-shared/fhe-macros}

\title{\textbf{Audit: fheCRDT Soundness and\\Empirical Performance}}
\author{
Automated Research Agent\\
\textit{Lux Industries}\\
\texttt{research@lux.network}
}
\date{2026-04-12}

\begin{document}
\maketitle
\luxcoverpage

\begin{abstract}
We audit the fheCRDT paper~\cite{fhecrdt-paper} and its reference
implementation in \texttt{github.com/luxfi/fhe}.  We verify lattice
properties (commutativity, associativity, idempotence) for all
implemented CRDT types and measure empirical merge throughput against
the paper's claims.  Key findings: (1) the LWW-Register implementation
is correct and preserves all three lattice properties under encryption;
(2) G-Counter, PN-Counter, G-Set, OR-Set, and RGA have \emph{no Go
implementation}---only the LWW-Register exists as executable code;
(3) measured throughput is \textbf{3--4 orders of magnitude below} the
paper's claims (0.24 LWW merges/sec vs.\ claimed 330/sec);
(4) the Solidity contract leaks timestamps in plaintext, undermining
the privacy model.
\end{abstract}

\section{Scope and Method}

\subsection{Materials Examined}

\begin{itemize}[noitemsep]
\item \texttt{cmd/crdt/main.go} --- LWW-Register demo (184 LOC)
\item \texttt{fhe.go}, \texttt{evaluator.go}, \texttt{encryptor.go},
      \texttt{decryptor.go} --- core FHE boolean gate engine
\item \texttt{examples/encrypted-crdt/contracts/EncryptedCRDT.sol}
      --- Solidity LWW-Register + OR-Set
\item \texttt{main.tex} --- the fheCRDT paper
\item \texttt{hanzo/base/crdt/} --- plaintext CRDT baseline
      (GCounter, PNCounter, LWWRegister, ORSet, MVRegister, RGA)
\end{itemize}

\subsection{Test Environment}

\begin{tabular}{ll}
\toprule
CPU         & Apple M1 Max (10-core, arm64) \\
RAM         & 64 GB \\
OS          & Darwin 25.4.0 \\
Go          & 1.26.0 \\
FHE params  & PN10QP27 (N=1024, Q=134215681, 128-bit security) \\
\bottomrule
\end{tabular}

\subsection{Method}

For each CRDT type claimed in the paper, we:
\begin{enumerate}[noitemsep]
\item Searched the codebase for a Go implementation.
\item Where code exists, wrote property-based tests for the three
      lattice laws plus observational equivalence to plaintext.
\item Ran Go benchmarks (\texttt{-benchtime=3s}) measuring wall-clock
      time per merge.
\item Compared against plaintext CRDT merge cost from
      \texttt{hanzo/base/crdt/}.
\end{enumerate}

\section{Soundness Findings}

\subsection{Implementation Coverage}

Table~\ref{tab:coverage} maps paper claims to actual code.

\begin{table}[h]
\centering
\caption{Paper claims vs.\ implementation status}
\label{tab:coverage}
\begin{tabular}{lccl}
\toprule
CRDT Type     & Paper & Go Code & Status \\
\midrule
G-Counter     & \ding{51} & \ding{55} & \textbf{Not implemented} \\
PN-Counter    & \ding{51} & \ding{55} & \textbf{Not implemented} \\
G-Set         & \ding{51} & \ding{55} & \textbf{Not implemented} \\
OR-Set        & \ding{51} & \ding{55} & \textbf{Not implemented} \\
LWW-Register  & \ding{51} & \ding{51} & Implemented (demo CLI) \\
RGA Sequence  & \ding{55} & \ding{55} & Not claimed, not implemented \\
\bottomrule
\end{tabular}
\end{table}

The paper's abstract claims ``practical implementations for counters,
sets, registers, and sequences.''  Only the LWW-Register exists as
runnable Go code.  The other four types are described algorithmically
in the paper but have zero implementation.

\subsection{LWW-Register: Lattice Properties}

We constructed encrypted registers with known plaintext values and
verified all three lattice laws empirically over FHE ciphertexts.

\begin{table}[h]
\centering
\caption{LWW-Register lattice property verification}
\label{tab:lattice}
\begin{tabular}{lccl}
\toprule
Property        & Test Cases & Result & Notes \\
\midrule
Commutativity   & (7,3)$\join$(12,5) & \textbf{PASS} & Both orders yield (12,5) \\
Idempotence     & (9,4)$\join$(9,4)  & \textbf{PASS} & Self-merge preserves state \\
Associativity   & 3-way merge        & \textbf{PASS} & (A$\join$B)$\join$C = A$\join$(B$\join$C) \\
\bottomrule
\end{tabular}
\end{table}

The merge algorithm (MSB-first comparator chain + MUX selection) is
correct.  After decryption, the encrypted merge result is
observationally equivalent to the plaintext LWW merge for all test
vectors.

\paragraph{Noise budget.}
Each MUX gate requires 3 bootstraps (AND, AND, OR).  The 4-bit
comparator requires $4 \times 2 = 8$ bootstraps (ANDNY + XNOR per
bit) plus $3 \times 2 = 6$ bootstraps for the reduction chain (AND +
OR + AND per non-MSB bit), totaling 14 comparator bootstraps.  The
value selection adds $4 \times 3 = 12$ MUX bootstraps and timestamp
selection adds another 12, for \textbf{38 total bootstraps per
4-bit LWW merge}.

Because each bootstrap resets noise to a fresh level (programmable
bootstrapping), there is \emph{no cumulative noise degradation} across
chained merges.  This is a strength of the TFHE/boolean-gate approach:
idempotence is preserved regardless of merge depth $k$.

\subsection{G-Counter: Theoretical Analysis}

The paper specifies G-Counter merge as per-replica $\FHEmax$, which
decomposes identically to the LWW timestamp comparator.  Since the
same comparator + MUX circuit is used, the lattice properties hold by
the same argument.  We verified this using \texttt{gcounterMerge} on
5 input pairs including equal values and extremes (0, 15).  All passed.

\subsection{G-Set: Theoretical Analysis}

G-Set merge is encrypted OR of membership flags.  OR is commutative,
associative, and idempotent over booleans.  Since bootstrapping
preserves bit correctness, encrypted OR inherits these properties.
We verified all 4 truth-table entries plus idempotence explicitly.
\textbf{PASS.}

\subsection{OR-Set: Gap Identified}

The paper's OR-Set description (Section 4.4) uses tag-set union for
merge, which is correct for the plaintext CRDT.  However, there is no
encrypted implementation.  The Solidity contract (\texttt{EncryptedCRDT.sol})
implements OR-Set with \emph{plaintext} tombstone flags and unique tags,
only encrypting the element \emph{values}.  The add/remove semantics leak
set membership patterns.

\subsection{Solidity Contract: Privacy Leak}

\texttt{EncryptedCRDT.sol} line 11: \texttt{uint256 timestamp} is
stored in plaintext.  The \texttt{mergeRegisters} function (line 88)
performs timestamp comparison in the clear:
\begin{lstlisting}[language=Solidity]
if (reg1.timestamp > reg2.timestamp) {
    merged = reg1.value;
\end{lstlisting}
This directly contradicts the paper's claim that ``merge operations
[are performed] homomorphically, keeping data encrypted at all
replicas.''  The timestamp---which reveals write ordering and
activity patterns---is fully exposed.

The Go implementation correctly compares timestamps homomorphically.
The Solidity contract does not.

\section{Empirical Benchmark Results}

\subsection{Measured Performance}

\begin{table}[h]
\centering
\caption{fheCRDT merge throughput (PN10QP27, M1 Max, Go 1.26)}
\label{tab:bench}
\begin{tabular}{lrrr}
\toprule
Operation & Latency & Ops/sec & Bootstraps \\
\midrule
Single bootstrap (Refresh)      & 108 ms  & 9.3     & 1 \\
G-Set merge (1 element)         & 114 ms  & 8.8     & 1 \\
G-Counter merge (4-bit, 1 slot) & 2,894 ms & 0.35   & $\sim$26 \\
LWW-Register merge (4-bit)      & 4,149 ms & 0.24   & $\sim$38 \\
Geo-convergence (50ms RTT)      & 4,146 ms & ---     & $\sim$38 \\
\bottomrule
\end{tabular}
\end{table}

\subsection{Comparison to Paper Claims}

\begin{table}[h]
\centering
\caption{Claimed vs.\ measured performance}
\label{tab:delta}
\begin{tabular}{lrrr}
\toprule
Metric & Paper Claim & Measured & Ratio \\
\midrule
G-Counter merge (10 replicas) & 500 ops/sec & 0.035 ops/sec$^\dagger$ & \textbf{14,286$\times$ slower} \\
LWW-Register merge            & 330 ops/sec & 0.24 ops/sec             & \textbf{1,375$\times$ slower} \\
G-Set merge (1K elements)     & 65 ops/sec  & 0.009 ops/sec$^\ddagger$ & \textbf{7,222$\times$ slower} \\
Convergence time              & 450 ms      & 4,146 ms                 & \textbf{9.2$\times$ slower} \\
LWW state size                & 64 bytes    & 136,456 bytes            & \textbf{2,132$\times$ larger} \\
\bottomrule
\end{tabular}

\footnotesize
$^\dagger$ Extrapolated: 10 replica slots $\times$ 2,894 ms each $= \sim$28.9 sec total.\\
$^\ddagger$ Extrapolated: 1000 elements $\times$ 114 ms each $= \sim$114 sec total.
\end{table}

\subsection{FHE vs.\ Plaintext Overhead}

\begin{table}[h]
\centering
\caption{Encrypted vs.\ plaintext merge latency}
\label{tab:overhead}
\begin{tabular}{lrrr}
\toprule
CRDT Type & Plaintext (ns) & Encrypted (ns) & Overhead \\
\midrule
G-Counter (1 slot) & 100       & 2,894,000,000     & $\sim$29 million$\times$ \\
LWW-Register       & 30        & 4,149,000,000     & $\sim$138 million$\times$ \\
G-Set (1 element)  & 326       & 114,000,000       & $\sim$350,000$\times$ \\
\bottomrule
\end{tabular}
\end{table}

\subsection{Ciphertext Size}

\begin{table}[h]
\centering
\caption{Ciphertext sizes (PN10QP27, N=1024)}
\label{tab:sizes}
\begin{tabular}{lr}
\toprule
Object & Size \\
\midrule
Single encrypted bit (serialized) & 17,057 bytes \\
Single encrypted bit (in-memory)  & 16,384 bytes \\
4-bit encrypted value             & 68,228 bytes \\
LWW-Register (4-bit val + 4-bit ts) & 136,456 bytes \\
\bottomrule
\end{tabular}
\end{table}

The paper claims 64 bytes for an LWW-Register state.  The measured
size is 136 KB---a 2,132$\times$ discrepancy.  The 64-byte figure
may refer to plaintext state size, not encrypted state.

\section{Bugs and Issues}

\begin{enumerate}
\item \textbf{Missing implementations} (\texttt{cmd/crdt/main.go}):
      Only LWW-Register is implemented.  The paper claims G-Counter,
      PN-Counter, G-Set, and OR-Set implementations exist.  They do not.

\item \textbf{Solidity timestamp leak}
      (\texttt{examples/encrypted-crdt/contracts/EncryptedCRDT.sol:11,57,102}):
      Timestamps are stored and compared in plaintext, breaking the
      homomorphic merge claim.

\item \textbf{Solidity OR-Set tombstone leak}
      (\texttt{EncryptedCRDT.sol:23}):
      The \texttt{removed} flag is a plaintext \texttt{bool}, leaking
      element removal events.

\item \textbf{Paper state size claim} (\texttt{main.tex} Table 4, line 372):
      ``LWW-Register: 64 bytes'' is incorrect for encrypted state.
      Actual encrypted size is 136 KB with PN10QP27 parameters.

\item \textbf{Paper throughput claims} (\texttt{main.tex} Table 4, line 367--373):
      All ops/sec figures are 3--4 orders of magnitude optimistic.
      No benchmark code or test environment is documented in the paper.
      The ``500 ops/sec'' G-Counter claim would require sub-2ms merge
      latency; a single bootstrap takes 108ms.

\item \textbf{Missing references in paper} (\texttt{main.tex:291,316}):
      References \texttt{[zap]} and \texttt{[luxcpp-fhe]} are cited but
      not defined in \texttt{references.bib}.

\item \textbf{No noise budget monitoring}: The implementation has no
      mechanism to track remaining noise budget or warn when ciphertexts
      approach decryption failure.  While TFHE bootstrapping resets noise
      per-gate, corrupted inputs from network transport or storage bit-flips
      would produce silent wrong answers.
\end{enumerate}

\section{Recommendation}

\textbf{fheCRDT is not production-ready for Hanzo Base as a sync backend.}

The gaps are:

\begin{enumerate}[noitemsep]
\item \textbf{4/5 claimed CRDT types have no implementation.}
      Only LWW-Register exists.  G-Counter, PN-Counter, G-Set, and OR-Set
      must be built before the framework can replace the plaintext CRDTs
      in \texttt{hanzo/base/crdt/}.

\item \textbf{Performance is 3--4 orders of magnitude too slow.}
      A 10-replica G-Counter merge would take $\sim$29 seconds.
      Base requires sub-100ms sync for interactive use.
      Closing this gap requires either:
      \begin{itemize}[noitemsep]
        \item GPU-accelerated bootstrapping (the paper cites 13--17$\times$
              speedup, still 2 orders of magnitude short), or
        \item CKKS/BFV packed ciphertexts for counter operations
              (avoiding per-bit bootstrapping), or
        \item Accepting that FHE merges happen asynchronously in batch,
              with plaintext CRDTs for real-time sync and periodic
              FHE reconciliation for privacy.
      \end{itemize}

\item \textbf{The Solidity contract contradicts the threat model.}
      Timestamps and tombstones are plaintext.  The contract needs
      encrypted Lamport clocks and encrypted set-membership flags
      with FHE comparison for merge.

\item \textbf{No delta-state support.}  The paper describes delta-state
      fheCRDTs (Section 5.1) but there is no implementation.  Full-state
      sync of 136 KB per LWW-Register per replica is prohibitive at scale.
\end{enumerate}

\paragraph{Path forward.}
The theoretical framework is sound---lattice properties are preserved
under FHE encryption, and the TFHE bootstrapping model avoids noise
accumulation.  The implementation needs:
(a) G-Counter and OR-Set Go implementations using the existing gate
primitives;
(b) CKKS-mode packed counters for practical throughput;
(c) delta-state protocol;
(d) the Solidity contract rewritten with encrypted timestamps.
Estimated effort: 4--6 engineer-weeks.

\section*{Acknowledgments}

Benchmarks ran on Apple M1 Max (10-core, 64 GB) under Darwin 25.4.0
with Go 1.26.0.  FHE parameters: PN10QP27 (N=1024, Q=134215681),
128-bit classical security level.

\begin{thebibliography}{9}
\bibitem{fhecrdt-paper}
Z.~Kelling, ``fheCRDT: Conflict-Free Replicated Data Types over
Encrypted State,'' Lux Industries, 2025.

\bibitem{shapiro2011}
M.~Shapiro, N.~Pregui\c{c}a, C.~Baquero, and M.~Zawirski,
``Conflict-free Replicated Data Types,'' SSS, 2011.

\bibitem{tfhe}
I.~Chillotti, N.~Gama, M.~Georgieva, and M.~Izabach\`ene,
``TFHE: Fast Fully Homomorphic Encryption over the Torus,''
ASIACRYPT, 2016.
\end{thebibliography}

\end{document}
